\documentclass[12pt,a4paper]{article}
\usepackage[utf8]{inputenc}
\usepackage[T2A]{fontenc}
\usepackage[russian]{babel}
\usepackage{geometry}
\usepackage{booktabs}
\usepackage{float}
\usepackage{amsmath}
\usepackage{hyperref}
\usepackage{listings}
\usepackage{xcolor}

\geometry{top=2cm, left=2cm, right=2cm, bottom=2cm}

\lstset{
  language=Python,
  basicstyle=\footnotesize\ttfamily,
  tabsize=2,
  breaklines=true,
  frame=single,
}

\title{Теория и вопросы защиты по всем четырём методам модуля 2:\\линейная регрессия, k-NN, дерево решений, логистическая регрессия}
\author{Ануфриев А.С. 465029}
\date{08.11.2026}

\begin{document}
\maketitle

\section{Введение}
В данном документе собран теоретический материал по четырём методам надзоруемого обучения, изученным в рамках модуля 2 (лабораторные работы №3--№6): линейная регрессия, метод k-ближайших соседей (k-NN), дерево решений и логистическая регрессия. Все методы были реализованы ``с нуля'' на языке Python с использованием только NumPy и Pandas без готовых библиотек машинного обучения (sklearn и аналоги). Документ предназначен для подготовки к защите: в конце каждого раздела приведены вероятные вопросы с ответами, а в завершении --- сравнительный обзор методов.

\section{Линейная регрессия (метод наименьших квадратов)}

\subsection{Суть метода}
Линейная множественная регрессия описывает зависимость результативного признака $y$ от $p$ факторов $x_1,\dots,x_p$ линейной моделью
\begin{equation}
y = a + b_1 x_1 + \dots + b_p x_p + \varepsilon ,
\end{equation}
где $a$ --- свободный член, $b_j$ --- регрессионные коэффициенты, $\varepsilon$ --- случайная ошибка. Эмпирическая оценка:
\begin{equation}
\hat{y} = a + b_1 x_1 + \dots + b_p x_p .
\end{equation}
В матричной записи $\hat{Y} = X B$, где в матрицу $X$ добавлен первый столбец единиц для свободного члена.

\subsection{Метод наименьших квадратов (МНК)}
Оценки параметров получаются минимизацией суммы квадратов остатков
\begin{equation}
Q(B) = \sum_i e_i^2 = (Y - XB)^\top (Y - XB).
\end{equation}
Дифференцирование по $B$ и приведение производной к нулю даёт нормальное уравнение
\begin{equation}
B = (X^\top X)^{-1} X^\top Y .
\end{equation}
Это аналитическое (замкнутое) решение: матричные операции одной итерации дают оптимальные коэффициенты при условии обратимости $X^\top X$.

\subsection{Качество модели}
Коэффициент детерминации
\begin{equation}
R^2 = 1 - \frac{\sum_i (y_i - \hat{y}_i)^2}{\sum_i (y_i - \bar{y})^2}
\end{equation}
показывает долю дисперсии целевой переменной, объясняемую моделью. $R^2 = 1$ --- идеальное описание, $R^2 = 0$ --- модель не точнее среднего значения.

\subsection{Предположения МНК (теорема Гаусса---Маркова)}
\begin{enumerate}
\item модель специфицирована верно (зависимость линейна);
\item между факторами нет строгой линейной зависимости (отсутствие мультиколлинеарности);
\item ошибки случайны, их математическое ожидание равно нулю, дисперсия постоянна (гомоскедастичность);
\item ошибки не коррелированы между собой.
\end{enumerate}
При выполнении этих условий оценки МНК являются наилучшими линейными несмещёнными оценками (BLUE).

\subsection{Особенности реализации в работе}
В лабораторной работе №3 коэффициенты вычислены собственной функцией на NumPy по формуле $B=(X^\top X)^{-1}X^\top Y$ без готовых средств регрессии. Признаки стандартизированы (вычет среднего, деление на стандартное отклонение), целевая переменная не нормировалась. Построено три модели: по доходу ($\approx 0.50$), по доходу и возрасту дома ($\approx 0.54$) и с синтетическим признаком $rooms\_per\_household = total\_rooms / households$ ($\approx 0.52$).

\subsection{Преимущества}
\begin{itemize}
\item простота реализации и интерпретации: каждый коэффициент показывает изменение цели при изменении на единицу соответствующего признака;
\item аналитическое решение, нет необходимости в итерационной оптимизации и подборе коэффициента обучения;
\item высокая скорость обучения и предсказания;
\item хорошая база для сравнения более сложных моделей.
\end{itemize}

\subsection{Недостатки и ограничения}
\begin{itemize}
\item предполагает линейную зависимость между признаками и целевой переменной;
\item чувствительна к мультиколлинеарности (в работе признаки $total\_rooms$, $households$ и их отношение создали строгую функциональную связь, что не улучшило модель 3);
\item чувствительна к выбросам, так как минимизирует сумму квадратов ошибок;
\item не подходит для классовых (дискретных) целей без модификаций.
\end{itemize}

\subsection{Вопросы защиты с ответами}
\begin{enumerate}
\item \textbf{Что такое метод наименьших квадратов и в чём его суть?}\\
МНК --- способ оценки параметров регрессионной модели, минимизирующий сумму квадратов отклонений наблюдаемых значений от предсказанных моделью. Аналитически решение задаётся формулой $B = (X^\top X)^{-1}X^\top Y$.
\item \textbf{Зачем нужен первый столбец единиц в матрице $X$?}\\
Он соответствует свободному члену $a$ --- значению $y$, когда все факторы равны нулю, и позволяет смещать регрессионную плоскость по оси $y$.
\item \textbf{Что означает коэффициент детерминации $R^2$?}\\
Доля дисперсии целевой переменной, объясняемая моделью. Чем ближе к единице, тем лучше. В работе лучшая модель --- $R^2 \approx 0.54$.
\item \textbf{Зачем нужно разделение на обучающую и тестовую выборки?}\\
Чтобы оценить качество модели на данных, не использовавшихся при обучении, и выявить переобучение. В работе --- разделение 80/20 (13 600 / 3 400 записей).
\item \textbf{Зачем нужна стандартизация признаков?}\\
Приводит признаки к сравнимому масштабу, упрощает интерпретацию коэффициентов и устойчивость вычислений.
\item \textbf{Что такое мультиколлинеарность и как она проявилась в работе?}\\
Строгая линейная зависимость между признаками. В третьей модели $total\_rooms$, $households$ и $rooms\_per\_household = total\_rooms/households$ линейно связаны, модель 3 не превзошла модель 2 ($R^2 \approx 0.52 < 0.54$).
\item \textbf{Какой признак сильнее всего влияет на стоимость жилья?}\\
Доход населения ($median\_income$) --- сильнейший отдельный фактор ($R^2 \approx 0.50$ в однофакторной модели).
\end{enumerate}

\section{Метод k-ближайших соседей (k-NN)}

\subsection{Суть метода}
k-NN --- метод надзоруемой классификации (и регрессии). Объект относится к классу, наиболее часто встречающемуся среди $k$ ближайших соседей в пространстве признаков. Расстояние обычно евклидово:
\begin{equation}
d(\mathbf{x}_i, \mathbf{x}_j) = \sqrt{\sum_{p=1}^{P} (x_{ip} - x_{jp})^2}.
\end{equation}
Класс определяется majoritaire голосованием среди $k$ ближайших обучающих примеров. Это ``ленивый'' алгоритм: явного этапа обучения нет, модель просто хранит обучающую выборку, а все вычисления выполняются на этапе предсказания.

\subsection{Выбор параметра $k$ и масштабирование}
Параметр $k$ определяет соседей, участвующих в голосовании. Малое $k$ делает модель чувствительной к шуму и выбросам (резкие границы классов), большое $k$ --- более устойчивой, но может сглаживать границы и стирать мелкие структуры. Оптимальное $k$ подбирается экспериментально (например, кросс-валидацией). Перед применением признаки масштабируются (стандартизация), иначе признаки с большими абсолютными значениями доминируют в расстоянии.

\subsection{Качество классификации}
Оценивается долей верных предсказаний (accuracy) и матрицей ошибок: строки --- истинные классы, столбцы --- предсказанные; диагональ --- верные ответы, внедиагональные элементы --- ошибки между конкретными классами.

\subsection{Особенности реализации в работе}
В лабораторной работе №4 k-NN реализован на NumPy: парные евклидовые расстояния вычисляются матрично через формулу $\|a-b\|^2=\|a\|^2+\|b\|^2-2a\cdot b$, затем сортировка, выбор $k$ ближайших и голосование через \texttt{np.bincount}. Готовые средства классификации не использовались. Построено две модели: со случайным набором признаков (5 из 13) и с фиксированным набором (6 признаков); исследованы $k=3,5,10,15$ на датасете о вине (178 записей, 3 класса, разделение 80/20). Лучшая точность --- 0.9444 (Модель 1 при $k=3,5$; Модель 2 при $k=10,15$). Основные ошибки --- путаница между классами 2 и 3.

\subsection{Преимущества}
\begin{itemize}
\item простота понимания и реализации;
\item нет этапа обучения (обучение --- просто запоминание данных);
\item непараметричность: не предполагает конкретной функциональной формы зависимости;
\item естественно работает с любым числом классов.
\end{itemize}

\subsection{Недостатки и ограничения}
\begin{itemize}
\item предсказание медленное: для каждого объекта нужны расстояния до всех обучающих примеров ($O(N \cdot P)$ на запрос);
\item высокая потребность в памяти (хранение всей обучающей выборки);
\item ``проклятие размерности'': при большом числе признаков расстояния теряют смысл;
\item чувствительность к выбору $k$, метрике расстояния и масштабу признаков;
\item плохо работает с категориальными признаками без предварительного кодирования.
\end{itemize}

\subsection{Вопросы защиты с ответами}
\begin{enumerate}
\item \textbf{К какому классу задач относится k-NN?}\\
Надзоруемая классификация (и регрессия), опирающаяся на расстояния в пространстве признаков.
\item \textbf{Как выбирается параметр $k$?}\\
Экспериментально (кросс-валидация). Малое $k$ --- чувствительность к шуму, большое $k$ --- размытие границ классов. В работе рассматривались $k = 3, 5, 10, 15$.
\item \textbf{Какая метрика расстояния использовалась и почему?}\\
Евклидово расстояние --- стандарт для k-NN на числовых признаках; адекватно измеряет близость после стандартизации.
\item \textbf{Зачем нужна стандартизация признаков перед k-NN?}\\
Признаки имеют разные размерности и диапазоны (например, Alcohol $\sim$ 13, Proline $\sim$ 747). Без масштабирования признаки с большими значениями доминируют в расстоянии.
\item \textbf{Чем отличаются Модель 1 и Модель 2 в работе?}\\
Модель 1 использует случайный набор признаков (5 из 13), Модель 2 --- фиксированный набор (6 признаков). Это оценивает влияние состава признаков на качество.
\item \textbf{Что показывает матрица ошибок?}\\
Распределение предсказаний по истинным и предсказанным классам; диагональ --- верные ответы, внедиагональ --- ошибки между классами.
\item \textbf{Почему ошибки приходятся на классы 2 и 3?}\\
Эти классы ближе друг к другу в пространстве признаков, чем класс 1, что затрудняет их раздел k-NN.
\end{enumerate}

\section{Деревья решений}

\subsection{Суть метода}
Дерево решений --- иерархическая структура правил классификации: каждый внутренний узел содержит проверку по какому-то признаку, ребра соответствуют возможным значениям признака, а листья содержат метки классов. Каждый объект попадает в один лист, давая покрывающий набор конъюнкций. Построение рекурсивно сверху вниз: в узле выбирается признак для разбиения, множество примеров делится на подмножества по значениям признака, для каждого создаётся потомок.

\subsection{Критерий разбиения (подход C4.5)}
Пусть $T$ --- множество примеров. Энтропия:
\begin{equation}
\mathrm{Info}(T) = -\sum_{j=1}^{k} \frac{freq(C_j,T)}{|T|} \log_2 \frac{freq(C_j,T)}{|T|}.
\end{equation}
После разбиения по признаку $X$ со значениями $A_1,\dots,A_n$ (подмножества $T_i$):
\begin{equation}
\mathrm{Info}_X(T) = \sum_{i=1}^{n} \frac{|T_i|}{|T|} \mathrm{Info}(T_i), \qquad
\mathrm{splitinfo}(X) = -\sum_{i=1}^{n} \frac{|T_i|}{|T|} \log_2 \frac{|T_i|}{|T|}.
\end{equation}
Нормированный прирост информации:
\begin{equation}
\mathrm{GainRatio}(X) = \frac{\mathrm{Info}(T) - \mathrm{Info}_X(T)}{\mathrm{splitinfo}(X)}.
\end{equation}
На каждом шаге выбирается признак с максимальным $\mathrm{GainRatio}$. Множитель $\mathrm{splitinfo}$ снижает преимущество признаков с большим числом уникальных значений.

\subsection{Особенности небинарного дерева}
В отличие от бинарных деревьев (CART), количество потомков узла не ограничено: каждое значение признака порождает отдельного потомка. В работе дерево реализовано через узел со ссылками на потомков (словарь ``значение признака $\rightarrow$ дочерний узел''), а не список-списков. Лист хранит класс большинства и долю объектов положительного класса, которая используется как вероятность предсказания.

\subsection{Условия остановки}
Рекурсия останавливается, если: все примеры в узле одного класса, достигнут предел максимальной глубины, признаки для разбиения исчерпаны, прирост информации равен нулю или недостаточно примеров.

\subsection{Особенности реализации в работе}
В лабораторной работе №5 построено небинарное дерево для датасета UCI Mushroom (грибы, съедобный / ядовитый, 22 номинативных признака). Пропуски (в признаке stalk-root) удалены; осталось 5644 записи (3488 съедобных, 2156 ядовитых). Отобрано случайно $\sqrt{22} \approx 4$ признака (зерно 42): cap-shape, cap-surface, gill-color, stalk-color-above-ring. Стратифицированное разделение 70/30 (3951 / 1693). Дерево глубины 4, 85 внутренних узлов, 63 листа. Метрики (положительный класс --- ядовитый): Accuracy 0.8547, precision 0.8692, recall 0.7295, AUC-ROC 0.9416, AUC-PR 0.9223. Все метрики и площади под кривыми вычислены собственным кодом (метод трапеций).

\subsection{Преимущества}
\begin{itemize}
\item высокая интерпретируемость: правила легко читаются и объясняются;
\item естественно работает с номинативными (категориальными) признаками;
\item Capturing нелинейных зависимостей и взаимодействий признаков без явной функциональной формы;
\item мало требований к предобработке (не нужно масштабирование).
\end{itemize}

\subsection{Недостатки и ограничения}
\begin{itemize}
\item склонность к переобучению, особенно на глубоких деревьях (требуется ограничение глубины, обрезка, выбор подмножества признаков);
\item нестабильность: небольшие изменения данных могут привести к другому дереву;
\item жадный алгоритм выбора признака не гарантирует глобально оптимальное дерево;
\item сложно моделировать аддитивные эффекты (например, сумму признаков) без глубоких веток.
\end{itemize}

\subsection{Вопросы защиты с ответами}
\begin{enumerate}
\item \textbf{Что представляет собой дерево решений?}\\
Иерархическая структура правил классификации; узлы --- проверки признаков, ребра --- значения признаков, листья --- классы.
\item \textbf{Чем алгоритм C4.5 отличается от бинарных деревьев?}\\
Количество потомков узла не ограничено, каждое значение признака порождает отдельного потомка.
\item \textbf{По какому критерию выбирается признак для разбиения?}\\
По максимальному нормированному приросту информации $\mathrm{GainRatio}(X)$.
\item \textbf{Как считается энтропия?}\\
$\mathrm{Info}(T) = -\sum_j p_j \log_2 p_j$, где $p_j$ --- доля примеров класса $C_j$.
\item \textbf{Что делает splitinfo?}\\
Корректирует прирост информации, снижая преимущество признаков с большим числом уникальных значений.
\item \textbf{Почему удалялись некоторые записи в данных о грибах?}\\
В признаке stalk-root встречались пропуски ``?'', такие объекты исключены перед обучением.
\item \textbf{Сколько признаков использовалось и как выбирались?}\\
$\sqrt{22} \approx 4$ признака, отобранных случайным образом с фиксированным зерном генератора.
\item \textbf{Как вычислялись Accuracy, precision, recall?}\\
По формулам из матрицы ошибок, собственным кодом без библиотек.
\item \textbf{Что такое AUC-ROC и AUC-PR?}\\
Площади под ROC- и PR-кривыми (пары FPR/TPR и Recall/Precision при разных порогах); площадь --- методом трапеций, без библиотек.
\end{enumerate}

\section{Логистическая регрессия}

\subsection{Суть метода}
Логистическая регрессия --- линейная модель для бинарной классификации. Модель оценивает вероятность принадлежности объекта к положительному классу:
\begin{equation}
P(y=1 \mid x) = \sigma(z), \qquad z = x^\top w,
\end{equation}
где $w$ --- вектор параметров (включая свободный член), $x$ --- вектор признаков, $\sigma(\cdot)$ --- сигмоидная функция
\begin{equation}
\sigma(z) = \frac{1}{1 + e^{-z}}, \qquad \sigma'(z) = \sigma(z)(1-\sigma(z)).
\end{equation}
Сигмоида отображает любое действительное число на отрезок $(0,1)$, что даёт интерпретацию выхода как вероятности. Порог классификации обычно равен 0.5: объект относится к классу 1, если $p \ge 0.5$.

\subsection{Функция потерь и обучение}
Используется функция логарифмических потерь (кросс-энтропия, log-loss):
\begin{equation}
L(w) = -\frac{1}{N}\sum_{i=1}^{N}\Bigl[y_i\ln p_i + (1-y_i)\ln(1-p_i)\Bigr],
\end{equation}
где $p_i = \sigma(x_i^\top w)$. Функция выпукла по $w$, поэтому градиентные методы сходятся к глобальному минимуму. Градиент:
\begin{equation}
\nabla_w L = \frac{1}{N}X^\top(p-y),
\end{equation}
обновление весов градиентным спуском:
\begin{equation}
w \leftarrow w - \eta \nabla_w L,
\end{equation}
где $\eta$ --- скорость обучения. Слишком большое $\eta$ может привести к дивергенции, слишком маленькое --- к медленной сходимости.

\subsection{Оптимизация Ньютона}
Вместо градиентного спуска используется второй порядок. Гессиан:
\begin{equation}
H = \frac{1}{N}X^\top W X, \qquad W = \mathrm{diag}(p_i(1-p_i)),
\end{equation}
обновление:
\begin{equation}
w \leftarrow w - H^{-1}\nabla_w L.
\end{equation}
Метод Ньютона обычно сходится за меньшее число итераций, но требует обращения матрицы гессиана размера $(d+1)\times(d+1)$ на каждой итерации, что дороже при большом числе признаков.

\subsection{Метрики классификации}
По матрице ошибок (TP, TN, FP, FN):
\begin{equation}
\text{Accuracy} = \frac{TP+TN}{TP+TN+FP+FN},\quad
\text{Precision} = \frac{TP}{TP+FP},\quad
\text{Recall} = \frac{TP}{TP+FN},\quad
\text{F1} = \frac{2\cdot\text{Precision}\cdot\text{Recall}}{\text{Precision}+\text{Recall}}.
\end{equation}
F1-Score --- гармоническое среднее precision и recall, устойчивое к дисбалансу классов.

\subsection{Особенности реализации в работе}
В лабораторной работе №6 реализована логистическая регрессия ``с нуля'' на NumPy/Pandas для датасета Титаника (891 запись, цель Survived, доля выживших 38,4\%). Предобработка без утечки данных: разделение 80/20 (712/179) до заполнения пропусков и стандартизации, статистики только из обучения, one-hot кодирование порта посадки, 9 признаков в модели. Исследованы гиперпараметры: learning rate $\{0.01, 0.1, 0.5\}$, итерации $\{100, 500, 1000\}$, методы градиентного спуска и Ньютона. Наилучший результат --- градиентный спуск, lr = 0.1, 500 итераций: accuracy 0.7765, precision 0.6866, recall 0.7077, F1 0.6970 (TP = 46, TN = 93, FP = 21, FN = 19). Ньютон сходится быстро, но не даёт преимущества в точности на тесте.

\subsection{Преимущества}
\begin{itemize}
\item вероятностный вывод (выдаёт оценку вероятности класса, а не только метку);
\item быстрая сходимость, простая реализация;
\item хорошая интерпретируемости: знаки и величины весов показывают влияние признаков;
\item хорошая базовая модель для бинарной классификации.
\end{itemize}

\subsection{Недостатки и ограничения}
\begin{itemize}
\item линейная граница разделения классов: $p \ge 0.5$ эквивалентно $x^\top w \ge 0$;
\item плохо справляется со сложными нелинейными зависимостями и взаимодействиями признаков без инженерии признаков;
\item чувствительна к дисбалансу классов (точность по меньшему классу может быть низкой при высоком общем accuracy);
\item требует масштабирования признаков для устойчивой сходимости градиентного спуска.
\end{itemize}

\subsection{Вопросы защиты с ответами}
\begin{enumerate}
\item \textbf{В чём отличие логистической регрессии от линейной?}\\
Линейная регрессия предсказывает непрерывное значение и минимизирует среднеквадратичную ошибку. Логистическая регрессия для классификации: через сигмоиду выдаёт вероятность в $(0,1)$, функцией потерь служит кросс-энтропия (log-loss), порог 0.5 превращает вероятность в класс 0 или 1.
\item \textbf{Зачем нужна сигмоида и какова её производная?}\\
Отображает любое $z$ на $(0,1)$, давая интерпретацию как вероятности. Производная $\sigma'(z) = \sigma(z)(1-\sigma(z))$ используется в градиенте и гессиане.
\item \textbf{Как выбирается порог классификации, и что изменится при его сдвиге?}\\
Стандартный порог --- 0.5. Повышение порога снижает число ложноположительных ответов (растёт precision) и увеличивает ложноотрицательные (падает recall). Понижение --- наоборот.
\item \textbf{Зачем нужна стандартизация признаков?}\\
Признаки имеют разные диапазоны (возраст до 80, стоимость билета до 512). Стандартизация приводит их к сравнимому масштабу, ускоряет и стабилизирует сходимость градиентного спуска.
\item \textbf{В чём преимущество и недостаток метода Ньютона?}\\
Преимущество --- обычно меньше итераций до сходимости благодаря второй производной. Недостаток --- необходимость обращать матрицу гессиана размера $(d+1)\times(d+1)$ на каждой итерации, что дороже при большом числе признаков.
\item \textbf{Можно ли сказать, что логистическая регрессия --- линейная модель?}\\
Да: решение принимается по линейной комбинации признаков $z = x^\top w$; поскольку сигмоида монотонна, $p \ge 0.5 \iff x^\top w \ge 0$ --- линейная граница в пространстве признаков.
\end{enumerate}

\section{Сравнительный обзор методов}

\subsection{Таблица сравнения}
\begin{table}[H]
\centering
\caption{Сравнительные характеристики четырёх методов модуля 2.}
\label{tab:summary}
\small
\begin{tabular}{p{2.4cm}p{2.1cm}p{2.6cm}p{3.4cm}p{3.2cm}}
\toprule
Метод & Тип задачи & Качество на тесте (это работа) & Преимущества & Ограничения \\
\midrule
Линейная регрессия (МНК) & Регрессия & Жильё Калифорнии, $R^2$ до 0.54 & Простота, интерпретируемость, аналитическое решение & Линейная зависимость, чувствительность к мультиколлинеарности и выбросам \\
k-NN & Классификация/регрессия & Вино, accuracy до 0.9444 & Нет обучения, непараметричность, любое число классов & Медленное предсказание, чувствительность к $k$, масштабу и размерности \\
Дерево решений & Классификация & Грибы, accuracy 0.8547, AUC-ROC 0.9416 & Интерпретируемость, номинативные признаки, нелинейность & Риск переобучения, нестабильность, жадность алгоритма \\
Логистическая регрессия & Бинарная классификация & Титаник, accuracy 0.7765, F1 0.6970 & Вероятности, быстрая сходимость, хорошая база & Линейная граница, сложности с нелинейными зависимостями \\
\bottomrule
\end{tabular}
\end{table}

\subsection{Когда какой метод наиболее эффективен}
\begin{itemize}
\item \textbf{Линейная регрессия} --- когда цель непрерывная и связь с признаками приближённо линейна (прогноз цен, спроса), нужна интерпретируемость эффектов факторов.
\item \textbf{k-NN} --- на маленьких/средних выборках, когда расстояние между объектами осмысленно, нет явной функциональной зависимости (рекомендательные системы, распознавание образов в низкоразмерном пространстве).
\item \textbf{Дерево решений} --- при смешанных числовых и категориальных признаках, когда важна понятная логика принятия решения (кредитный скоринг, медицинская диагностика).
\item \textbf{Логистическая регрессия} --- для бинарной классификации с вероятностным выводом (отток, скоринг рисков, выживание), особенно как базовая, легко обучаемая модель.
\end{itemize}

\section{Общие вопросы защиты по модулю}
\begin{enumerate}
\item \textbf{В чём разница между регрессией и классификацией?}\\
Регрессия предсказывает непрерывное числовое значение, классификация --- дискретный метку класса.
\item \textbf{Зачем нужны обучающая и тестовая выборки?}\\
Для оценки обобщающей способности модели на данных, не использовавшихся при обучении, и выявления переобучения.
\item \textbf{Что такое переобучение и как с ним бороться?}\\
Модель слишком точно подстраивается под обучающие данные (включая шум) и плохо обобщается. Борьба: ограничение сложности модели, кросс-валидация, выбор подмножества признаков, регуляризация.
\item \textbf{Почему важна стандартизация признаков?}\\
Методы на основе расстояний (k-NN) и градиентного спуска (логистическая регрессия) чувствительны к масштабу признаков; стандартизация приводит их к сравнимым диапазонам.
\item \textbf{Почему все методы в работе реализованы без sklearn?}\\
Для понимания внутреннего устройства алгоритмов: формул, оптимизации, расчёта метрик от первого лица.
\item \textbf{Как выбиралась ``главная'' метрика в каждой работе?}\\
Для регрессии --- $R^2$; для классификации --- accuracy, precision, recall, F1 (для несбалансированных классов F1 более информативна, чем accuracy); для деревьев --- дополнительно AUC-ROC и AUC-PR.
\end{enumerate}

\end{document}